\section*{第 \ref{ch:g7:removing-wastes} 章——排出废物}

\begin{solution}{exo:g7:removing-wastes:1}
一种含氮废物，是\omterm{def:g5:first-look-at-cells:cell}{细胞}在利用
构建物质的\omterm{def:g7:digestion-and-nutrients:families}{营养物质}时留下的。它在流动中无害，
但若大量堆积便会中毒，因此必须
持续排出。
\end{solution}

\begin{solution}{exo:g7:removing-wastes:2}
肺：\omterm{def:g7:what-is-respiration:respiration}{二氧化碳}。\omterm{def:g7:removing-wastes:kidneys}{肾脏}：\omterm{prop:g7:removing-wastes:rubbish}{尿素}，以及过剩的
水分和盐类。\omterm{def:g1:the-five-senses:sense}{皮肤}：热量，以及\omterm{rem:g7:removing-wastes:sweat}{汗液}中的水分和盐类
（并含少量\omterm{prop:g7:removing-wastes:rubbish}{尿素}）。
\end{solution}

\begin{solution}{exo:g7:removing-wastes:3}
两个，拳头大小、形似豆子，位于腰背偏下处，
各由一条粗大的\omterm{def:g7:heart-and-circulation:vessels}{动脉}供血。\omterm{def:g7:removing-wastes:kidneys}{尿液}从每侧
经输尿管流入\omterm{def:g7:removing-wastes:kidneys}{膀胱}，\omterm{def:g7:removing-wastes:kidneys}{膀胱}是可随时排空的储存器官。
\end{solution}

\begin{solution}{exo:g7:removing-wastes:4}
先大量过滤：每天有几十升\omterm{prop:g4:journey-of-food:absorb}{血液}中的水分
连同其中有用与有害的物质被压过滤器。
随后沿弯曲的肾小管节俭地重吸收，
概括来说：过滤几十升，
约一升半以\omterm{def:g7:removing-wastes:kidneys}{尿液}形式排出。
\end{solution}

\begin{solution}{exo:g7:removing-wastes:5}
全部被重吸收的是\omterm{prop:g7:organs-at-work:bill}{葡萄糖}（以及几乎全部的水）。
按需调节重吸收量的是水分和盐类，视当天需要而定。
完全不被重吸收的是\omterm{prop:g7:removing-wastes:rubbish}{尿素}——这正是本题的关键。
\end{solution}

\begin{solution}{exo:g7:removing-wastes:6}
重吸收量会随情况调节：大量饮水后，重吸收的水分减少，
因而尿量多、颜色浅；出汗多或饮水少时，
几乎全部水分都被重吸收，因而尿量少、颜色深。
\omterm{def:g7:removing-wastes:kidneys}{肾脏}借此调控体内的水分收支。
\end{solution}

\begin{solution}{exo:g7:removing-wastes:7}
\omterm{def:g7:what-is-respiration:respiration}{二氧化碳}是气体，气体在\omterm{def:g7:lungs-and-gas-exchange:alveoli}{肺泡}这一空气交界处穿过。
\omterm{prop:g7:removing-wastes:rubbish}{尿素}不是气体，无法随\omterm{def:g7:what-is-respiration:respiration}{呼吸}排出，
而是\omterm{def:g7:oxygen-in-the-water:dissolved}{溶解}在\omterm{prop:g4:journey-of-food:absorb}{血液}中流动，并在液体交界处，
即“筛子”处，被很好地滤出。每种废物都经由
与其形态相\omterm{def:g4:adapted-to-their-world:adaptation}{适应}的途径排出：分工，是化学性质
与结构的精妙配合。
\end{solution}

\begin{solution}{exo:g7:removing-wastes:8}
滤器让\omterm{prop:g7:organs-at-work:bill}{葡萄糖}自由通过，这是正常的；
之后的重吸收阶段理应将其全部回收。
若\omterm{prop:g7:organs-at-work:bill}{葡萄糖}出现在\omterm{def:g7:removing-wastes:kidneys}{尿液}中，说明重吸收
不足或已受损：要么血糖过高，
重吸收跟不上，要么重吸收功能本身失常——
两种情况都值得就医关注。
\end{solution}

\begin{solution}{exo:g7:removing-wastes:9}
\omterm{def:g3:bones-and-muscles:muscle}{肌肉}\omterm{def:g7:heart-and-circulation:vessels}{毛细血管}——\omterm{def:g7:heart-and-circulation:vessels}{静脉}——右\omterm{def:g7:heart-and-circulation:rooms}{心房}和右\omterm{def:g7:heart-and-circulation:rooms}{心室}——
肺循环（\omterm{prop:g7:removing-wastes:rubbish}{尿素}随之流动，肺无法处理它）——
左心——通往\omterm{def:g7:removing-wastes:kidneys}{肾脏}的体循环\omterm{def:g7:heart-and-circulation:vessels}{动脉}——
在滤过单位处被过滤——\emph{未}被重吸收——
沿集合管向下，最终流入\omterm{def:g7:removing-wastes:kidneys}{膀胱}。
\end{solution}

\begin{solution}{exo:g7:removing-wastes:10}
进入途径：不经分解，在\omterm{def:g4:journey-of-food:tract}{小肠}\omterm{prop:g7:digestion-and-nutrients:absorption}{绒毛}处进入
\omterm{prop:g4:journey-of-food:absorb}{血液}。运输：先经\omterm{prop:g7:digestion-and-nutrients:absorption}{肝脏}，再经两条循环。
用途：作为每个\omterm{def:g5:first-look-at-cells:cell}{细胞}湿润化学反应的介质，
并为散热系统供水。排出途径：由\omterm{def:g7:removing-wastes:kidneys}{尿液}（经重吸收调节）、
\omterm{rem:g7:removing-wastes:sweat}{汗液}（运动时占较大份额）、呼出水汽，
以及少量随残余物排出。
\end{solution}

\begin{solution}{exo:g7:removing-wastes:11}
滤器借助精细的膜从\omterm{prop:g4:journey-of-food:absorb}{血液}中滤去废物；
机器即体外滤过装置。这类治疗需要反复进行，
因为\omterm{def:g5:first-look-at-cells:cell}{细胞}不断产生\omterm{prop:g7:removing-wastes:rubbish}{尿素}
（见\cref{prop:g7:removing-wastes:rubbish}）：每周运行两次的滤过装置，
必须清除数日的积累，
而\omterm{def:g7:removing-wastes:kidneys}{肾脏}的滤器每分钟都在工作。
\end{solution}

\begin{solution}{exo:g7:removing-wastes:12}
举例来说：同一\omterm{prop:g4:journey-of-food:absorb}{血液}循肺循环和体循环流动，由四室\omterm{def:g4:heart-and-blood:heart}{心脏}驱动，
在\omterm{def:g7:lungs-and-gas-exchange:alveoli}{肺泡}处装载\omterm{def:g7:what-is-respiration:respiration}{氧气}，在\omterm{prop:g7:digestion-and-nutrients:absorption}{绒毛}处装载\omterm{def:g7:digestion-and-nutrients:families}{营养物质}，
并将两者输送到每个\omterm{def:g5:first-look-at-cells:cell}{细胞}的\omterm{def:g7:heart-and-circulation:vessels}{毛细血管}中，
在那里\omterm{def:g7:what-is-respiration:respiration}{呼吸作用}这把缓慢的“火”借\omterm{def:g7:what-is-respiration:respiration}{氧气}燃烧燃料，
为一切活动提供动力。同一\omterm{prop:g4:journey-of-food:absorb}{血液}又带走这把“火”的灰烬——
\omterm{def:g7:what-is-respiration:respiration}{二氧化碳}送至\omterm{def:g7:lungs-and-gas-exchange:alveoli}{肺泡}，\omterm{prop:g7:removing-wastes:rubbish}{尿素}和多余物质送至\omterm{def:g7:removing-wastes:kidneys}{肾脏}的
滤过与重吸收装置，因此这四道交界共同服务于
同一套运作体系。每个器官都是这一循环中的一站，
而这一循环的存在，是为了让\omterm{def:g5:first-look-at-cells:cell}{细胞}中的“火”
始终不缺燃料、\omterm{def:g7:what-is-respiration:respiration}{氧气}，
也不缺废物的清运。
\end{solution}

\begin{solution}{pb:g7:removing-wastes:1}
\textbf{1.} 在\omterm{def:g4:journey-of-food:tract}{小肠}的\omterm{prop:g7:digestion-and-nutrients:absorption}{绒毛}处穿过，
进入\omterm{def:g7:heart-and-circulation:vessels}{毛细血管}，由\omterm{prop:g4:journey-of-food:absorb}{血液}网络接收。

\textbf{2.} \omterm{prop:g7:digestion-and-nutrients:absorption}{肝脏}。路径的安排服务于与\omterm{def:g7:digestion-and-nutrients:families}{营养物质}
共享同一\omterm{prop:g4:journey-of-food:absorb}{血液}的物质——分类与储备——同时
水分则只是沿\omterm{def:g3:sorting-living-things:criterion}{标准}路线经过\omterm{prop:g7:digestion-and-nutrients:absorption}{肝脏}。

\textbf{3.} 肺循环和体循环，
通过四室\omterm{def:g4:heart-and-blood:heart}{心脏}严格交替进行。

\textbf{4.} 它是\omterm{def:g5:first-look-at-cells:cell}{细胞}化学的介质——
\omterm{def:g4:journey-of-food:digestion}{消化液}与燃烧物质皆以\omterm{def:g7:oxygen-in-the-water:dissolved}{溶解}态运行——并充当
\omterm{prop:g4:journey-of-food:absorb}{血液}所运输一切物质的载体；它还
为下午以\omterm{rem:g7:removing-wastes:sweat}{汗液}形式消耗的散热储备提供来源。

\textbf{5.} \omterm{def:g1:the-five-senses:sense}{皮肤}：\omterm{rem:g7:removing-wastes:sweat}{汗液}——水分与盐分——其
当天真正的作用是为像熔炉般灼热的工作
\omterm{def:g3:bones-and-muscles:muscle}{肌肉}降温（见\cref{rem:g7:organs-at-work:warmth}）。

\textbf{6.} 出汗消耗了水分储备；
\omterm{def:g7:removing-wastes:kidneys}{肾脏}的重吸收阶段几乎将滤过的水分全部
收回，\omterm{def:g7:removing-wastes:kidneys}{尿液}少而浓、颜色深——
而口渴则要求补足这部分储备。

\textbf{7.} \omterm{def:g7:what-is-respiration:respiration}{二氧化碳}：\omterm{def:g3:bones-and-muscles:muscle}{肌肉}\omterm{def:g7:heart-and-circulation:vessels}{毛细血管}、\omterm{def:g7:heart-and-circulation:vessels}{静脉}、右心
\omterm{def:g7:lungs-and-gas-exchange:alveoli}{肺泡}——呼出。\omterm{prop:g7:removing-wastes:rubbish}{尿素}：同样的回流途径，
随后经过肺、左心、肾\omterm{def:g7:heart-and-circulation:vessels}{动脉}，被滤过
并送往\omterm{def:g7:removing-wastes:kidneys}{膀胱}。

\textbf{8.} \omterm{prop:g7:removing-wastes:rubbish}{尿素}以及经\omterm{rem:g7:removing-wastes:sweat}{汗液}作用后\omterm{prop:g4:journey-of-food:absorb}{血液}中的盐分：
废物的产生从不停歇，因此滤过器官也不能停歇。
\emph{尿量}减少——重吸收留住了水分
但清除过程仍在继续。

\textbf{9.} \omterm{def:g7:lungs-and-gas-exchange:alveoli}{肺泡}表面湿润，每一次
\omterm{def:g7:what-is-respiration:respiration}{呼吸}都把变暖、湿润的空气排出——即镜面上雾气
的水分。任何玻璃的雾气，
都有一部分以\omterm{def:g4:breathing:movements}{呼气}中的\omterm{def:g6:exploring-our-environment:conditions}{湿度}形式离开身体。

\textbf{10.} 进入：\omterm{prop:g7:digestion-and-nutrients:absorption}{绒毛}。储存与利用：\omterm{prop:g4:journey-of-food:absorb}{血液}
体积、\omterm{def:g5:first-look-at-cells:cell}{细胞}化学、散热储备。排出：\omterm{def:g7:removing-wastes:kidneys}{尿液}
（经重吸收调节）、\omterm{rem:g7:removing-wastes:sweat}{汗液}（比赛所占份额）、呼出
的水分，以及少量未被吸收、随
残渣排出的部分。

\textbf{11.} 由于每个器官都从同一条\omterm{prop:g4:journey-of-food:absorb}{血液}中取得供应，
并向其中排放废物。一个衰竭的肺会使每块\omterm{def:g3:bones-and-muscles:muscle}{肌肉}缺乏供给，
衰竭的肾会毒害每个\omterm{def:g5:first-look-at-cells:cell}{细胞}，
硬化的\omterm{def:g7:heart-and-circulation:vessels}{动脉}会对其下游的任一站点
形成节流。一个网络意味着一种命运，
逐站如此。

\textbf{12.} \omterm{def:g7:what-is-respiration:respiration}{呼吸作用}——每个\omterm{def:g5:first-look-at-cells:cell}{细胞}中缓慢的“燃烧”：
进入、运输、燃烧、构建，以及每一条
废物通道，都是为了供养它、供给它，
并带走它的灰烬。
\end{solution}
